\documentclass[10pt,a4paper]{amsart}
\usepackage[margin=19mm]{geometry}
\usepackage{amsmath,amssymb,mathtools}
\usepackage[scr=rsfs,cal=euler]{mathalfa}
\usepackage{xfrac}
\usepackage[unicode,hidelinks,bookmarks=false]{hyperref}
\newcommand{\ie}{i.e.\ }
\newcommand{\cf}{cf.\ }
\newcommand{\resp}{respectively\ }
\newcommand{\Subsection}{\subsection}
\providecommand{\nobreakdash}{\nobreak\hbox{-}}
\setlength{\emergencystretch}{2em}
\multlinegap0pt
\usepackage{bm}
\usepackage{bbm}
\usepackage{dsfont}
\usepackage{xspace}
\usepackage{cancel}
\usepackage{mathtools}
\usepackage{amsthm}
\makeatletter
\@ifundefined{theorem}{\newtheorem{theorem}{Theorem}}{}
\@ifundefined{example}{\newtheorem{example}[theorem]{Example}}{}
\@ifundefined{proposition}{\newtheorem{proposition}[theorem]{Proposition}}{}
\makeatother
\def\R{\mathbb R}
\def\cone{\mathop{\mathrm{cone}}}
\def\sK{\mathcal K}
\newcommand{\mb}[1]{{\textbf {\textit#1}}}
\renewcommand{\ge}{\geqslant}
\renewcommand{\le}{\leqslant}
\title[Exponential actions defined by vector configurations, Gale d]{Exponential actions defined by vector configurations, Gale duality, and moment-angle manifolds}
\author{Taras Panov}
\date{Source: arXiv:2411.03366v2 (CC BY 4.0)}
\begin{document}
\maketitle

\noindent\emph{Source and licence.} Exponential actions defined by vector configurations, Gale duality, and moment-angle manifolds. Version arXiv:2411.03366v2 (CC BY 4.0), distributed under the Creative Commons Attribution 4.0 International licence, as recorded in the arXiv licence metadata for this exact version and shown on the abstract page https://arxiv.org/abs/2411.03366v2. Full rights notice: licenses/arxiv-2411.03366-CC-BY-4.0.txt.\par\smallskip

\noindent\textit{Editorial scope.} Counted unit: the Theorem whose source label is \texttt{comcom} in the article (source0.tex lines 820--823, source label \texttt{comcom}), delivered together with the article's own complete original proof of it (source0.tex lines 824--917, one \texttt{proof} environment of 94 lines). No mathematics is written, repaired or re-derived: statement and proof are the article's own text line for line. Required context retained uncounted: affui (lines 475--477); galeindspan (lines 510--513); 1dimfan (lines 795--809). Every definition, display and label of those lines is retained unchanged. Omitted from this package and not redistributed: 1--474, 478--509, 514--794, 810--819, 918--3247. Nothing inside a retained excerpt is omitted or altered: every retained body is delivered exactly as the article's lines are written, so no omission receipt of any kind accompanies this package. Citations: the retained excerpts carry no citation command at all, so no bibliography entry of the article is transcribed and no reference is redistributed. Reference adaptation: none was needed and none was made; every reference inside the delivered unit resolves inside it, so no unresolved reference or citation is produced. Printed slips kept literal: no printed word, symbol, hypothesis or number of the counted statement or of its proof was altered. Layout and class: the article's own class is \texttt{amsart} with packages including amssymb, amscd, color, graphicx, xy, tikz, tikz-cd, tikz-3dplot, hyperref; the wrapper is the lane's pinned \texttt{amsart} editorial wrapper at 10pt on A4, which restates the theorem-like environments the retained text needs and adds the article's own macro definitions that the retained text needs (\texttt{\textbackslash R}, \texttt{\textbackslash cone}, \texttt{\textbackslash ge}, \texttt{\textbackslash le}, \texttt{\textbackslash mb}, \texttt{\textbackslash sK}), with extra wrapper packages (bm, bbm, dsfont, xspace, cancel, mathtools). The delivered numbering is the unit's own. Assets: no retained span contains a figure environment, a graphics-inclusion command, a PDF-page inclusion, a table or a TikZ picture; no image file is copied into this package, the package declares no asset dependency and no image file is redistributed with it. Licence: version arXiv:2411.03366v2 of this article is distributed under the Creative Commons Attribution 4.0 International licence, as recorded in the arXiv licence metadata for this exact version and shown on the abstract page https://arxiv.org/abs/2411.03366v2. A full rights notice is delivered at licenses/arxiv-2411.03366-CC-BY-4.0.txt. Characters: every retained excerpt is pure ASCII, so the delivered engine, XeLaTeX with its default Latin Modern text font, reports no missing character.
\begin{equation}\label{affui}
  U_I=\{(x_1,\ldots,x_m)\in\R^m\colon x_j\ne0\text{ for }j\notin I\}.
\end{equation}

\begin{proposition}\label{galeindspan}
%Let $\Gamma$ and $\mathrm{A}$ be Gale dual configurations of $m$ vectors in $V^*$ and $W^*$, %respectively. 
For any $I\subset[m]$, the vectors in $\mathrm{A}_I=\{\mb a_i\colon i\in I\}$ are linearly independent in $W^*$ if and only if $\Gamma_{\widehat I}=\{\gamma_j\colon j\notin I\}$ spans~$V^*$.  
\end{proposition}

\begin{example}\label{1dimfan}
Consider the action of $V=\R$ on $\R^2$ given by 
\[
  (v,(x_1,x_2))\mapsto(e^{v}x_1,e^{v}x_2).
\]
We have $\Gamma=(\gamma_1,\gamma_2)=(1,1)$, $\mathrm A=(\mb a_1,\mb a_2)=(1,-1)$. Let $\sK=\sK(\Gamma)=\{\varnothing,\{1\},\{2\}\}$, so that $U(\sK)=\R^2\setminus\{\mathbf 0\}$.
The vectors $\mb a_1,\mb a_2$ form a one-dimensional fan in~$\R$, so the above action of $\R$ on $U(\sK)=\R^2\setminus\{\mathbf 0\}$ is proper and the quotient is homeomorphic to a circle (a smooth manifold).

Now consider the action of $V=\R$ on $\R^2$ given by 
\[
  (v,(x_1,x_2))\mapsto(e^{v}x_1,e^{-v}x_2).
\]
This time we have $\Gamma=(\gamma_1,\gamma_2)=(1,-1)$, $\mathrm A=(\mb a_1,\mb a_2)=(1,1)$.
The relative interiors of the cones generated by $\mb a_1$ and $\mb a_2$ have a non-empty intersection (they coincide), so the action of $\R$ on $U(\sK)=\R^2\setminus\{\mathbf 0\}$ is not proper and the quotient is a non-Hausdorff space.
\end{example}

\begin{theorem}\label{comcom}
Assume that $\{\mathcal K,\mathrm A\}$ is a triangulated configuration defining a simplicial fan~$\Sigma$.
The quotient smooth manifold $U(\sK)/V$ is compact if and only if the fan $\Sigma$ is complete.
\end{theorem}

\begin{proof}
Assume that $\Sigma$ is a complete fan.
Because $U(\sK)/V$ is metrisable, it is enough to prove that any sequence of points in $U(\sK)/V$ contains a convergent subsequence.

Take a sequence of points in $U(\sK)/V$ and let $\{\mb x^{(n)},n=1,2,\ldots\}$ be its lift to the total space, $\mb x^{(n)}=(x_1^{(n)} , \dots , x_m^{(n)}) \in U(\sK)$ . We shall change the sequence $\{\mb x^{(n)}\}$ to $\{\widetilde{\mb x}^{(n)}\}$ so that $\mb x^{(n)}$ and $\widetilde{\mb x}^{(n)}$ lie in the same $V$-orbit for each $n=1,2,\ldots$ and $\{\widetilde{\mb x}^{(n)}\}$ has a convergent  subsequence  in $U(\sK)$. Its projection will be the required subsequence in $U(\sK)/V$. 

By passing to a subsequence we may assume that  all elements of the sequence $\{\mb x^{(n)}\}$ have the same sets of zero coordinates. That is, there is $J\in\sK$ such that $x^{(n)}_j=0$ for $j\in J$ and $x^{(n)}_j\ne0$ for $j\notin J$, for all~$n$.

We argue by induction on $\dim W^*=m-k$. Let $\dim W^*=1$.  There is only one complete one-dimensional fan (corresponding to $\sK=\{\varnothing,\{1\},\{2\}\}$), so the theorem holds in this case (see Example~\ref{1dimfan}).

If $J\ne\varnothing$, then $(\mb v\cdot\mb x^{(n)})_j=0$ for any $j\in J$, $\mb v\in V$ and $n=1,2,\ldots$ Consider the link subcomplex $\mathop\mathrm{link}_{\sK}J=\{I\in\sK\colon I\cup J\in\sK,\,I\cap J=\varnothing\}$. Then the data 
$\{\mathop\mathrm{link}_{\sK}J,\mb a_j\colon j\notin J\}$ define a complete fan in the quotient space $W^*/\langle\mb a_j\colon j\in J\rangle$, which has dimension less than~$W^*$. Therefore, we can apply the induction assumption in this case.

We therefore may assume $J=\varnothing$, that is, $x_j^{(n)}\ne0$ for any $n$ and $j\in[m]$.

We denote $\R^\times=\R\setminus\{0\}$. Given $\mb x=(x_1,\ldots,x_m)\in(\R^\times)^m$, define 
\[
  \ell(\mb x)=\sum_{i=1}^m\mb a_i\log|x_i|\in W^*.
\]
%We think of $\ell(\mb x)$ as a linear function on $W$, whose value on $\mb w$ is 
%$\sum_{i=1}^m\langle\mb a_i,\mb w\rangle\log|x_i|$. 
For any $\mb v\in V$ we have
\begin{multline}\label{lconst}
\ell(\mb v\cdot\mb x)=\sum_{i=1}^m\mb a_i\log\bigl(e^{\langle\gamma_i,\mb v\rangle}|x_i|\bigr)
\\=\sum_{i=1}^m\mb a_i\langle\gamma_i,\mb v\rangle+\sum_{i=1}^m\mb a_i\log|x_i|=\sum_{i=1}^m\mb a_i\log|x_i|=\ell(\mb x),
\end{multline}
where the second-to-last equality follows by Gale duality.

Now we get back to the sequence $\{\mb x^{(n)}\}$ and consider $-\ell(\mb x^{(n)})\in W^*$. As $\Sigma$ is a complete fan in $W^*$, we can find $I\in\sK$ with $|I|=\dim W^*$ such that infinitely many terms of the sequence $\{-\ell(\mb x^{(n)})\}$ lie in $\cone\mathrm A_I$. Hence, by passing to a subsequence, we may assume that
$\{-\ell(\mb x^{(n)})\}\in\cone\mathrm A_I$ for $n=1,2,\ldots$.

By Proposition~\ref{galeindspan}, the vector configuration $\Gamma_{\widehat I}$ is a basis of $V^*$, so we can find $\mb v^{(n)}\in V$ satisfying
$$
   \langle \gamma_i ,\mb v^{(n)}\rangle = -\log|x_{i}^{(n)}|,\quad i \notin I.
$$
and set
\[
  \widetilde{\mb x}^{(n)}=\mb v^{(n)}\cdot{\mb x}^{(n)}.
\]
We have
\[
  |\widetilde x_i^{(n)}| = |e^{\langle \gamma_i,\mb v^{(n)}\rangle} x_i^{(n)}|=1,\quad i\notin I.
\]
This together with~\eqref{lconst} implies
\[
-\sum\limits_{i \in I}\mb a_i \log |\widetilde x_i^{(n)}| =-\ell(\widetilde{\mb x}^{(n)})=  -\ell(\mb x^{(n)}) \in \cone\mathrm A_I.
\]
It follows that $\log |\widetilde x_i^{(n)}|\le0$ for any $i\in[m]$ and $n=1,2,\ldots$. Therefore, 
$\{\widetilde{\mb x}^{(n)}\}$ is a bounded sequence, so its subsequence has a limit it $\R^m$. As $|\widetilde x_i^{(n)}|=1$ for $i\notin I$, the limit may have zero coordinates $x_i$ only for $i\in I$, so the limit belongs to $U(\sK)$, as needed.

\smallskip

Now assume that $U(\sK)/V$ is compact. Take $\mb a\in W^*$. To prove that $\Sigma$ is complete, we need to show that $\mb a$ belongs to a cone of~$\Sigma$. Write 
\begin{equation}\label{aalpha}
  \mb a=\sum_{i=1}^m\alpha_i\mb a_i, \quad \alpha_i\in\R.
\end{equation}
For any $x>0$ consider
\[
  (x^{\alpha_1},\ldots,x^{\alpha_m})\in (\R^\times)^m\subset U(\sK)
\]
and let $[x^{\alpha_1},\ldots,x^{\alpha_m}]$ denote the image in $U(\sK)/V$. Setting $x=x^{(n)}$, $n=1,2,\ldots$, where $x^{(n)}>0$ converges to zero as $n\to\infty$ and using the fact that $U(\sK)/V$ is compact, we observe that there exists a limit
%$[x^{\alpha_1},\ldots,x^{\alpha_m}]$ must have a limit in $U(\sK)/V$ as $x\to0$. Let
\[
  \lim_{x\to0}[x^{\alpha_1},\ldots,x^{\alpha_m}]=[y_1,\ldots,y_m]\in U(\sK)/V.
\]
Let 
\[
  I=\{i\in[m]\colon y_i=0\}\in\sK.
\]
This is well-defined as the set of zero coordinates of a point in $U(\sK)$ is constant along its $V$-orbit. Furthermore, $\mb y=(y_1,\ldots,y_m)\in U_I$, where $U_I\in U(\sK)$ is given by~\eqref{affui}.

For any $\mb w\in W$ and $\mb x= (x_1,\ldots,x_m)\in (\R^\times)^m$, define
\[
  f_{\mb w}(x_1,\ldots,x_m)=\prod_{i=1}^m|x_i|^{\langle\mb a_i,\mb w\rangle}.
\]
Then $f_{\mb w}$ defines a function $(\R^\times)^m/V\to\R$, because $f_{\mb w}(\mb x)=e^{\langle\ell(\mb x),\mb w\rangle}$ for $\mb x\in(\R^\times)^m$ and $\ell(\mb x)$ is constant on $V$-orbits (see above). Consider the dual cone
\[
  (\cone \mathrm A_I)^\mathsf{v}=
  \{\mb w\in W\colon\langle\mb a_i,\mb w\rangle\ge 0\text{ for }i\in I\}.
\]
If $\mb w\in(\cone \mathrm A_I)^\mathsf{v}$, then $f_{\mb w}$ extends to a continuous  function on~$U_I/V$.

Now $\lim_{x\to0}[x^{\alpha_1},\ldots,x^{\alpha_m}]=[\mb y]$ in $U_I/V$, implying that 
\[
  \lim_{x\to0}f_{\mb w}(x^{\alpha_1},\ldots,x^{\alpha_m})=f_{\mb w}(\mb y)
\]  
 for $\mb w\in(\cone \mathrm A_I)^\mathsf{v}$. On the other hand,
\[
  \lim_{x\to0}f_{\mb w}(x^{\alpha_1},\ldots,x^{\alpha_m})  
  =\lim_{x\to 0} x^{\langle\sum_{i=1}^m\alpha_i\mb a_i,\mb w\rangle}
  =\lim_{x\to 0} x^{\langle\mb a,\mb w\rangle}.
\]
The latter limit exists if and only if $\langle\mb a,\mb w\rangle\ge0$. Hence, $\langle\mb a,\mb w\rangle\ge0$ for any $\mb w\in(\cone \mathrm A_I)^\mathsf{v}$. This implies  that $\mb a\in\cone \mathrm A_I$, and therefore the fan $\Sigma$ is complete.
\end{proof}

\end{document}
